\chapter*{Declaration of GenAI Usage}

[The use of generative artificial intelligence (GenAI) must be declared in accordance with SDU regulations and guidance. Students are responsible for their use of AI-generated output, and the intellectual work presented in the report must remain their own. The declaration should make it clear to the reader how generative AI contributed to the final work.

If generative AI tools were used during the project, briefly describe:
\begin{itemize}
    \item Tools used, e.g. Microsoft Copilot, ChatGPT, GitHub Copilot, or other AI-based tools
    \item Purpose and extent of use, e.g. brainstorming, structuring the report, literature search, language improvement, summarization, analysis, generation or modification of text, programming, or generation of other material
    \item Where relevant, indicate the use of generative AI at the appropriate place in the report, source code, or other project material
\end{itemize}

If generative AI was not used, state this explicitly.

Students must follow the current SDU rules and guidance concerning generative AI. Consult:
\begin{itemize}
    \item SDU Library: SDU regulations and guidance on generative AI
    \item Software Engineering S\o{}nderborg: Guide to the use of generative AI
\end{itemize}

The Software Engineering guidance provides more detailed instructions for different forms of AI use. For example, AI-generated or AI-paraphrased text requires appropriate declaration/citation, while uses such as report structuring, brainstorming, language assistance, analysis, and programming should be described at a relevant point in the assignment.]
